% chapters/ch02_memory_toolchain.tex — Chapter 2: Memory & Toolchain
% Scope (PLAN.md): memory map of an embedded program, stack vs heap,
% .text/.data/.bss, linker scripts, startup code (what runs before main),
% compilation stages, Makefiles, stack overflow & memory corruption diagnosis.
% Budget 30 pp, mix 15 Q&A / 12 code / 3 concept.
% All host-compilable snippets verified with gcc -Wall -Wextra -std=c11 (zero
% warnings). Target-only fragments are marked and excluded from host builds.

\chapter{Memory \& Toolchain}
\label{ch:memory-toolchain}

\begin{partnote}
\textbf{Scope} --- The memory map of an embedded program $\rightarrow$ stack
vs heap $\rightarrow$ \code{.text}/\code{.data}/\code{.bss}/\code{.rodata}
$\rightarrow$ the build pipeline (preprocess $\rightarrow$ compile $\rightarrow$
assemble $\rightarrow$ link) $\rightarrow$ linker scripts (VMA/LMA)
$\rightarrow$ startup code (what runs before \code{main}) $\rightarrow$
Makefiles $\rightarrow$ diagnosing stack overflow and memory corruption.
This is the chapter that explains \emph{where} the Chapter 1 keywords
actually live: where a \code{static} ends up, why a \code{const} can sit in
flash, and what zero-initialises your globals before \code{main}.
\end{partnote}

If Chapter 1 was the language, this is the machine the language runs on. An
embedded interviewer asks these questions to find out whether you have ever
debugged a board that ``runs in the debugger but not from reset'' --- the
answer to which lives in the linker script and the startup code, not in your C.

%=====================================================================
\section{Primer}
%=====================================================================

\subsection{The memory map of an embedded program}

A linked firmware image is divided into \emph{sections}, each with a purpose
and a home:
\begin{itemize}
\item \code{.text} --- executable code and (usually) string literals. On an
  MCU it lives in \textbf{flash} and runs in place.
\item \code{.rodata} --- read-only data: \code{const} objects, literal tables,
  lookup tables. Also in flash.
\item \code{.data} --- initialised read/write globals/statics (e.g.\
  \code{int g = 5;}). The \emph{values} are stored in flash; the
  \emph{variables} live in RAM; startup copies flash $\rightarrow$ RAM.
\item \code{.bss} --- zero-initialised read/write globals/statics (e.g.\
  \code{int g;} or \code{static int s;}). Occupies RAM but \textbf{not} flash;
  startup zeroes it.
\item \textbf{stack} --- automatic (local) variables, return addresses, saved
  registers; grows (on ARM/x86) from a high address downward.
\item \textbf{heap} --- \code{malloc} territory; grows upward toward the stack.
  Often unused or tightly bounded in flight software.
\end{itemize}
The single most useful mental model: \emph{flash holds the program and the
initial values; RAM holds everything that changes; startup code bridges the
two before \code{main} runs.}

\subsection{The build pipeline}

A C file becomes a running image in four stages:
\begin{enumerate}
\item \textbf{Preprocess} (\code{cpp}): expand \code{\#include}, \code{\#define},
  and conditionals $\rightarrow$ a single translation unit of pure C.
\item \textbf{Compile} (\code{cc1}): translation unit $\rightarrow$ assembly
  for the target architecture.
\item \textbf{Assemble} (\code{as}): assembly $\rightarrow$ an object file
  (\code{.o}) of machine code plus a symbol table and unresolved references.
\item \textbf{Link} (\code{ld}): combine object files and libraries, resolve
  symbols, place sections at addresses per the \textbf{linker script}
  $\rightarrow$ the final \code{.elf}. A separate step (\code{objcopy})
  produces the raw \code{.bin}/\code{.hex} that gets flashed.
\end{enumerate}
Knowing which stage emits which error saves hours: a missing semicolon is a
\emph{compile} error; ``undefined reference to \code{foo}'' is a \emph{link}
error; ``\code{region FLASH overflowed}'' is a \emph{linker-script} error.

\subsection{Startup code and the linker script}

Between reset and \code{main} runs the \textbf{startup code} (often
\code{startup.s} + a C \code{Reset\_Handler}). On a Cortex-M, the reset vector
points at it; it (1) sets the stack pointer, (2) copies \code{.data} from its
load address in flash (LMA) to its run address in RAM (VMA), (3) zeroes
\code{.bss}, (4) runs C++ constructors / \code{init\_array}, then (5) calls
\code{main}. The \textbf{linker script} (\code{.ld}) declares the memory
regions (FLASH, RAM) and assigns sections to them, exporting the symbols
(\code{\_etext}, \code{\_sdata}, \code{\_edata}, \code{\_sbss}, \code{\_ebss})
that the startup code loops over. ``Works in debugger, fails from reset'' is
almost always a startup/linker-script bug: the debugger loaded RAM for you;
reset did not.

%=====================================================================
\section{Q\&A Bank}
%=====================================================================

\tierbanner{Tier 1 --- Screening (phone / HR-technical)}%
{Two to four sentences. Segment and pipeline basics --- a stumble here signals
no real board experience.}

\question{Q2.1}{Name the main memory segments of a C program and what each
holds.}
\code{.text} holds code (and often literals) in flash; \code{.rodata} holds
\code{const} read-only data in flash; \code{.data} holds initialised
read/write globals in RAM; \code{.bss} holds zero-initialised globals in RAM;
the \textbf{stack} holds locals and call frames; the \textbf{heap} serves
\code{malloc}. Flash keeps what's fixed, RAM keeps what changes.

\question{Q2.2}{Stack vs heap in one breath?}
The stack is automatic, fast, LIFO storage for locals and return addresses,
allocated and freed by the compiler as functions enter and exit; the heap is
manually managed (\code{malloc}/\code{free}), flexible but slower,
fragmentation-prone, and non-deterministic. On embedded the stack is small and
fixed; the heap is often avoided entirely.

\question{Q2.3}{What is the difference between \code{.data} and \code{.bss}?}
\code{.data} is for globals/statics with a non-zero initialiser --- their
initial values must be stored in flash and copied to RAM at startup.
\code{.bss} is for globals/statics that are zero (or uninitialised) --- they
take RAM but no flash, because startup just zeroes the region. So a big
zero-filled array costs RAM but not image size.

\question{Q2.4}{What are the four stages of building a C program?}
Preprocessing (expand includes/macros), compilation (C $\rightarrow$
assembly), assembly (assembly $\rightarrow$ object file), and linking (combine
objects, resolve symbols, place sections $\rightarrow$ executable). On embedded
a final \code{objcopy} turns the ELF into a flashable \code{.bin}/\code{.hex}.

\question{Q2.5}{Where do a global variable and a local variable live?}
A global (or any \code{static}) has static storage duration: it lives in
\code{.data} or \code{.bss} for the whole program. A local has automatic
storage duration: it lives on the stack and exists only while its function
runs --- which is why returning the address of a local is a classic dangling
pointer.

\tierbanner{Tier 2 --- Core technical (panel round)}%
{A short paragraph. The panel is checking you understand the path from reset
to \code{main} and how to read the toolchain's output.}

\question{Q2.6}{What runs before \code{main} on a bare-metal target?}
The startup code, reached via the reset vector. It initialises the stack
pointer, copies the \code{.data} section from its flash load address to its
RAM run address, zeroes the \code{.bss} section, runs any static
constructors / \code{init\_array} entries, and then calls \code{main}. If the
\code{.data} copy or \code{.bss} zeroing is missing or misconfigured, globals
hold garbage and the firmware misbehaves in ways that vanish under a debugger
(which loads RAM for you).

\question{Q2.7}{What is a linker script and what does it define?}
It is the file (\code{.ld}) that tells the linker the target's memory map ---
the address and size of each region (FLASH, RAM) --- and which output
sections go where, in what order. It assigns each section a run address (VMA)
and a load address (LMA), and exports boundary symbols (\code{\_sdata},
\code{\_edata}, \code{\_sbss}, \code{\_ebss}, \code{\_etext}) that the startup
code uses to copy \code{.data} and zero \code{.bss}. Get the region sizes
wrong and you get ``region overflowed''; get the LMA/VMA wrong and globals
don't initialise.

\question{Q2.8}{Explain VMA vs LMA with \code{.data}.}
The VMA (virtual/run address) is where a section is \emph{used} at run time;
the LMA (load address) is where its initial contents are \emph{stored} in the
image. For code in flash they're the same. For \code{.data} they differ: the
variables run from RAM (VMA in RAM) but their initial values are stored in
flash (LMA in flash), and startup copies LMA $\rightarrow$ VMA. That copy is
exactly why an initialised global has the right value the instant \code{main}
starts.

\question{Q2.9}{Why is the heap discouraged in embedded/flight software ---
from a memory-map view?}
The heap grows toward the stack in a fixed RAM region, so a runaway allocation
or fragmentation can collide with the stack (heap--stack collision) with no
MMU to catch it. \code{malloc} timing is non-deterministic (breaks WCET
analysis), fragmentation can fail an allocation despite enough total free
memory, and there is rarely a safe in-flight recovery. The standard answer:
allocate statically or from fixed-size pools at init, never in the control
loop --- which is what the Coding Gym pool/arena allocators below implement.

\question{Q2.10}{Declaration vs definition --- and what causes ``undefined
reference'' vs ``multiple definition''?}
A \emph{declaration} introduces a name and type (\code{extern int x;}); a
\emph{definition} also allocates storage (\code{int x;}). ``Undefined
reference'' is a \emph{link} error: something used a symbol that no object
file defines (forgot to compile/link the \code{.c}, or only declared it).
``Multiple definition'' is the opposite: two translation units each defined
the same external symbol (e.g.\ a non-\code{static} global or a function body
in a header). The fixes: put the single definition in one \code{.c} and
\code{extern} declarations in the header; mark file-local helpers
\code{static}; mark header inline functions \code{static inline}.

\question{Q2.11}{How would you find out how much flash and RAM your firmware
uses?}
Two standard tools. \code{size firmware.elf} prints the \code{text}/\code{data}/
\code{bss} totals --- flash usage is roughly \code{text + data}, RAM usage is
roughly \code{data + bss} (plus runtime stack/heap). The linker \textbf{map
file} (\code{-Wl,-Map=fw.map}) breaks it down per symbol and per section, so
you can see which object or array is eating memory and confirm sections landed
in the right regions. For static stack analysis you add tooling or the painting
technique in Q3.14.

\tierbanner{Tier 3 --- Trap \& depth (principal-engineer level)}%
{A paragraph plus the trap. These separate ``I flash binaries'' from
``I understand the image.''}

\question{Q3.12}{A global reads as 0 before \code{main} in one build and as
garbage in another. Explain.}
Zero-initialised globals live in \code{.bss}, which the C standard guarantees
is zero \emph{at program startup} --- but on bare metal that guarantee is only
as good as the startup code that zeroes the \code{.bss} region. If the startup
loop (or the linker symbols \code{\_sbss}/\code{\_ebss} it uses) is wrong or
absent, \code{.bss} holds whatever was in RAM at power-on --- garbage --- and
the bug appears only when running from reset, not under a debugger that
cleared RAM. \trap the candidate who says ``globals are always zero'' has never
brought up a board. The defensible answer names \code{.bss}, the startup
zeroing loop, and the debugger-vs-reset discrepancy --- the same class of
issue as the BEL ``works in the lab rig, fails on the unit'' defects.

\question{Q3.13}{Differentiate a stack overflow from heap corruption: cause,
symptom, and detection on a target without an MMU.}
A \emph{stack overflow} is the stack growing past its allotted RAM --- from
deep recursion, large local arrays, VLAs, or an ISR storm --- overwriting
adjacent memory (often \code{.bss}/\code{.data} or the heap), so unrelated
globals mysteriously change. \emph{Heap corruption} is writing outside an
allocated block or double-/use-after-free, which smashes allocator metadata and
typically faults on a later \code{malloc}/\code{free}, far from the real bug.
Detection without an MMU: a \textbf{stack canary/guard word} at the stack
limit (trip if it changes), \textbf{painting} the stack with a known pattern
and measuring the high-water mark, an \textbf{MPU} region marking the stack
guard no-access (faults on overflow), and \textbf{red-zone guard bytes} around
heap blocks. \trap juniors say ``it crashes.'' Seniors name the \emph{distance}
between corruption and symptom --- the reason these bugs are hard --- and give
a concrete detection mechanism for each. (Both are built in the Coding Gym.)

\question{Q3.14}{How do you measure worst-case stack usage on a deployed
embedded system?}
Two complementary methods. \textbf{Static}: compiler stack-usage output
(\code{-fstack-usage}) plus call-graph analysis gives a worst-case bound
\emph{except} through function pointers and recursion, which must be bounded by
hand. \textbf{Runtime ``painting''}: fill the whole stack region with a known
sentinel (e.g.\ \code{0xCD}) at startup, run the worst-case workload, then scan
from the far end for the first untouched sentinel --- the distance is the
high-water mark, and the remaining unpainted margin is your headroom. \trap the
trap is trusting a single test run; you must drive the genuine worst case
(deepest call path \emph{with} the highest-priority ISR nested on top) before
the number means anything. RTOSes expose this as ``stack high-water mark'' per
task --- implemented exactly as the painting problem below.

\question{Q3.15}{Walk the linker through resolving \code{int counter;} defined
in two \code{.c} files vs declared \code{extern} in a header --- why one links
and one doesn't.}
\code{extern int counter;} in a shared header is a \emph{declaration}: every
file that includes it knows the symbol's type but allocates nothing, and the
linker binds them all to the one real definition. \code{int counter;} placed
in the header instead becomes a \emph{tentative definition} in every including
file; with the traditional ``common symbol'' model the linker merged them, but
modern toolchains (GCC 10+, \code{-fno-common} by default) treat them as
multiple definitions and reject the link --- which is the correct, safe
behaviour. \trap the depth here is tentative definitions and \code{-fno-common}:
the historically-``working'' header-with-a-bare-global pattern is now a link
error, and the fix is the disciplined \code{extern}-in-header /
definition-in-one-\code{.c} split. Mark truly file-local state \code{static} so
it never participates in linking at all.

%=====================================================================
\section{Coding Gym}
%=====================================================================

\begin{partnote}
Memory/toolchain concepts are best learned by building the structures the
linker and allocator would otherwise hide. Every solution compiles clean under
\code{gcc -Wall -Wextra -std=c11}; the allocators and guards are the real
patterns used in flight-grade firmware that bans \code{malloc}.
\end{partnote}

%---------------------------------------------------------------
\begin{gymproblem}{Map the segments --- which variable lives where}

\textbf{Problem.} Write a program that takes the addresses of an initialised
global, an uninitialised global, a \code{const}, a function, a local, and a
heap block, and prints them. Predict their relative ordering and name each
segment.

\textbf{Hints.}
\begin{itemize}
\item \code{\&g\_init} is \code{.data}; \code{\&g\_uninit} is \code{.bss};
  \code{\&g\_const} is \code{.rodata}; \code{main} is \code{.text}.
\item Stack addresses are typically high; heap addresses sit above the static
  data. Exact values vary per platform.
\end{itemize}

\attempt

\textbf{Solution.}
\begin{cgym}
#include <stdio.h>
#include <stdlib.h>

int g_init = 5;                 /* .data   : initialised global  */
int g_uninit;                   /* .bss    : zero-init global     */
const int g_const = 7;          /* .rodata : read-only constant   */
int counter(void) { static int s; return ++s; }  /* static local -> .bss */

int main(void) {
    int local = 1;                          /* stack */
    int *heap = malloc(sizeof *heap);       /* heap  */
    if (!heap) return 1;
    *heap = 2;
    printf("data@%p bss@%p ro@%p text@%p stack@%p heap@%p\n",
           (void*)&g_init, (void*)&g_uninit, (void*)&g_const,
           (void*)(void(*)(void))main, (void*)&local, (void*)heap);
    printf("vals: %d %d %d %d %d %d\n",
           g_init, g_uninit, g_const, counter(), local, *heap);
    free(heap);
    return 0;
}
\end{cgym}

The \emph{values} are deterministic (\code{5 0 7 1 1 2}) --- note
\code{g\_uninit} and the \code{static s} both start at 0 because they're in
\code{.bss}. The \emph{addresses} vary by platform, but the ordering is
instructive: code and \code{.rodata} lowest (flash on an MCU), then
\code{.data}/\code{.bss}, then the heap growing up, and the stack high and
growing down.

\followups
\begin{enumerate}
\item \emph{``On an MCU, which of these are in flash vs RAM?''} $\rightarrow$
  \code{.text}/\code{.rodata}/the \code{.data} \emph{initialisers} in flash;
  \code{.data}/\code{.bss}/stack/heap in RAM.
\item \emph{``Why is \code{g\_uninit} zero without you writing 0?''}
  $\rightarrow$ \code{.bss} is zeroed by startup code before \code{main}.
\item \emph{``Move \code{g\_uninit} into \code{.data} --- how?''} $\rightarrow$
  Give it a non-zero initialiser; now its value costs flash.
\end{enumerate}
\end{gymproblem}

%---------------------------------------------------------------
\begin{gymproblem}{Which way does the stack grow?}

\textbf{Problem.} Determine at runtime whether the stack grows toward lower or
higher addresses.

\textbf{Hints.}
\begin{itemize}
\item Compare the address of a local in a callee against one from the caller.
\item A lower callee address means the stack grows downward.
\end{itemize}

\attempt

\textbf{Solution.}
\begin{cgym}
#include <stdio.h>

static int grows_down(int *parent_local) {
    int child_local;
    return (&child_local < parent_local);  /* callee frame below caller? */
}

int main(void) {
    int top;
    printf("stack grows %s\n", grows_down(&top) ? "down" : "up");
    return 0;
}
\end{cgym}

On ARM and x86 this prints ``down'': each call places its frame at a lower
address. (Strictly, comparing pointers into different frames is outside what
the standard defines, so this is a practical probe, not portable contract ---
say so if asked.) Why it matters: stack-overflow guards and MPU stack-limit
regions are placed at the \emph{low} end on a downward stack, and a deep call
chain or big local array marches toward that limit.

\followups
\begin{enumerate}
\item \emph{``Where do you put the stack guard, given the direction?''}
  $\rightarrow$ At the lowest stack address (the limit it grows toward).
\item \emph{``Why is this technically UB?''} $\rightarrow$ Relational
  comparison of pointers into different objects/frames is unspecified; use a
  documented stack-limit symbol on a real target.
\item \emph{``How does an ARM Cortex-M signal the stack limit in hardware?''}
  $\rightarrow$ \code{PSPLIM}/\code{MSPLIM} registers (ARMv8-M) or an MPU
  guard region.
\end{enumerate}
\end{gymproblem}

%---------------------------------------------------------------
\begin{gymproblem}{Fixed-block pool allocator (the \code{malloc} replacement)}

\textbf{Problem.} Implement \code{pool\_alloc}/\code{pool\_free} over a static
array of fixed-size blocks --- deterministic, no fragmentation.

\textbf{Hints.}
\begin{itemize}
\item A parallel \code{used[]} flag array tracks occupancy.
\item Allocation returns the first free block; free clears its flag.
\end{itemize}

\attempt

\textbf{Solution.}
\begin{cgym}
#include <stdio.h>
#include <stdint.h>
#include <stddef.h>

#define BLOCK_SIZE 32
#define NUM_BLOCKS 4
static uint8_t pool[NUM_BLOCKS][BLOCK_SIZE];
static uint8_t used[NUM_BLOCKS];

void *pool_alloc(void) {
    for (int i = 0; i < NUM_BLOCKS; i++)
        if (!used[i]) { used[i] = 1; return pool[i]; }
    return NULL;                       /* pool exhausted -> caller handles it */
}
void pool_free(void *p) {
    for (int i = 0; i < NUM_BLOCKS; i++)
        if (p == pool[i]) { used[i] = 0; return; }
}

int main(void) {
    void *a = pool_alloc(), *b = pool_alloc();
    pool_free(a);
    void *c = pool_alloc();            /* reuses a's block */
    printf("reused=%d b=%d\n", c == a, b != NULL);   /* reused=1 b=1 */
    return 0;
}
\end{cgym}

Every allocation is the same size and either succeeds or returns \code{NULL} ---
no fragmentation, bounded time, fully static RAM. This is the canonical pattern
for message buffers, CAN frames, or task control blocks in a system that bans
\code{malloc}. The \code{free}-list version (problem 8) makes alloc/free O(1).

\followups
\begin{enumerate}
\item \emph{``Make allocation O(1).''} $\rightarrow$ Keep a free list of block
  indices/pointers instead of scanning \code{used[]} --- see problem 8.
\item \emph{``How do you make it ISR-safe?''} $\rightarrow$ Brief critical
  section (disable interrupts) around the flag update, or a lock-free index
  stack.
\item \emph{``How do you detect a double free here?''} $\rightarrow$ Check the
  flag is set before clearing; flag a fault if freeing an already-free block.
\end{enumerate}
\end{gymproblem}

%---------------------------------------------------------------
\begin{gymproblem}{Bump / arena allocator with alignment}

\textbf{Problem.} Implement \code{arena\_alloc(size, align)} that hands out
aligned chunks from a static buffer by advancing an offset.

\textbf{Hints.}
\begin{itemize}
\item Round the current offset up to \code{align} before carving.
\item \code{align\_up(n, a) = (n + (a-1)) \& \textasciitilde(a-1)} for
  power-of-two \code{a}.
\end{itemize}

\attempt

\textbf{Solution.}
\begin{cgym}
#include <stdio.h>
#include <stdint.h>
#include <stddef.h>

static uint8_t arena[256];
static size_t  a_off = 0;

static size_t align_up(size_t n, size_t a) {
    return (n + (a - 1)) & ~(a - 1);     /* a must be a power of two */
}
void *arena_alloc(size_t size, size_t align) {
    size_t cur = align_up(a_off, align);
    if (cur + size > sizeof arena) return NULL;   /* out of arena */
    a_off = cur + size;
    return &arena[cur];
}

int main(void) {
    char     *p = arena_alloc(10, 1);
    uint32_t *q = arena_alloc(sizeof(uint32_t), 4);
    printf("p=%d q_aligned=%d\n", p != NULL, ((uintptr_t)q % 4u) == 0);
    return 0;
}
\end{cgym}

A bump allocator is the fastest possible allocator --- one add per allocation
--- and never fragments \emph{during} a run because it only moves forward. The
catch is you cannot free individual objects; you reclaim the whole region at
once (problem 12). Alignment matters because an unaligned \code{uint32\_t}
access faults or is slow on some cores --- the same alignment discipline behind
DMA buffers.

\followups
\begin{enumerate}
\item \emph{``Why must \code{align} be a power of two?''} $\rightarrow$ The
  \code{\& \textasciitilde(a-1)} mask only rounds correctly for powers of two.
\item \emph{``How do you free?''} $\rightarrow$ You don't free individually;
  reset the offset (region/arena reset) --- problem 12.
\item \emph{``Align a DMA buffer to 32 bytes --- change what?''} $\rightarrow$
  Pass \code{align = 32}; also ensure the arena itself is suitably aligned.
\end{enumerate}
\end{gymproblem}

%---------------------------------------------------------------
\begin{gymproblem}{Pointer alignment check before a word access}

\textbf{Problem.} Write \code{is\_aligned(p, a)} returning true when pointer
\code{p} meets alignment \code{a} (a power of two).

\textbf{Hints.}
\begin{itemize}
\item Cast to \code{uintptr\_t} and mask the low bits.
\item \code{(addr \& (a-1)) == 0} means aligned.
\end{itemize}

\attempt

\textbf{Solution.}
\begin{cgym}
#include <stdio.h>
#include <stdint.h>
#include <stddef.h>

static int is_aligned(const void *p, size_t a) {
    return ((uintptr_t)p & (a - 1)) == 0;   /* a must be a power of two */
}

int main(void) {
    uint32_t x = 0;
    uint8_t  buf[8] = {0};
    printf("x=%d misaligned=%d\n",
           is_aligned(&x, 4),        /* a uint32_t is 4-aligned -> 1 */
           is_aligned(buf + 1, 4));  /* odd offset -> 0              */
    return 0;
}
\end{cgym}

\code{uintptr\_t} is the integer type guaranteed to hold a pointer, so the
mask is well defined. Alignment checks guard the fast path: before a
word-at-a-time \code{memcpy}, before casting a byte buffer to a
\code{uint32\_t*}, and before handing a buffer to DMA --- on a Cortex-M0 or
older ARM an unaligned 32-bit access is a fault, not a slowdown.

\followups
\begin{enumerate}
\item \emph{``Round a pointer up to alignment.''} $\rightarrow$
  \code{(uint8\_t*)(((uintptr\_t)p + (a-1)) \& \textasciitilde(a-1))}.
\item \emph{``Why \code{uintptr\_t} not \code{int}?''} $\rightarrow$ An
  \code{int} may be narrower than a pointer; the cast would lose bits.
\item \emph{``Where does unaligned access actually fault?''} $\rightarrow$
  Cortex-M0/M0+/M1 and classic ARM for 32-bit loads; M3/M4 allow it for normal
  loads but not for \code{LDM}/\code{STM} or strongly-ordered memory.
\end{enumerate}
\end{gymproblem}

%---------------------------------------------------------------
\begin{gymproblem}{Red-zone guard bytes --- detect a buffer overrun}

\textbf{Problem.} Wrap a buffer in known guard words and provide a check that
reports whether anything wrote past its bounds.

\textbf{Hints.}
\begin{itemize}
\item Put a sentinel before and after the data inside one struct (layout is
  then deterministic).
\item Verify both sentinels are intact.
\end{itemize}

\attempt

\textbf{Solution.}
\begin{cgym}
#include <stdio.h>
#include <string.h>
#include <stdint.h>

#define GUARD 0xA5A5A5A5u
typedef struct {
    uint32_t front;        /* red zone before the data */
    uint8_t  data[16];
    uint32_t back;         /* red zone after the data  */
} guarded_t;

static void guarded_init(guarded_t *g) { g->front = GUARD; g->back = GUARD; }
static int  guarded_ok(const guarded_t *g) {
    return g->front == GUARD && g->back == GUARD;
}

int main(void) {
    guarded_t g;
    guarded_init(&g);
    memcpy(g.data, "safe", 5);          /* within bounds */
    printf("intact=%d\n", guarded_ok(&g));   /* intact=1 */
    return 0;
}
\end{cgym}

A write of more than 16 bytes into \code{data} would clobber \code{back}, and
\code{guarded\_ok} would return false --- turning silent corruption into a
detectable event you can assert on. Embedding the guards in a struct makes the
layout deterministic (unlike a loose canary next to a stack array, which the
compiler may reorder). This is the debugging-allocator / ``electric fence''
idea in miniature.

\followups
\begin{enumerate}
\item \emph{``How would a real heap debugger use this?''} $\rightarrow$ Pad
  every allocation with guard bytes and check them on \code{free}.
\item \emph{``Why guards inside a struct, not separate variables?''}
  $\rightarrow$ The compiler can reorder/realign separate locals; struct member
  order is fixed.
\item \emph{``What's the runtime cost, and when do you compile it out?''}
  $\rightarrow$ Extra RAM + check time; gate behind a \code{DEBUG} build.
\end{enumerate}
\end{gymproblem}

%---------------------------------------------------------------
\begin{gymproblem}{Model the startup code --- copy \code{.data}, zero
\code{.bss}}

\textbf{Problem.} Using arrays to stand in for the linker-defined regions,
write the two loops the startup code runs before \code{main}: copy the
\code{.data} initialisers from ``flash'' to ``RAM'' and zero the \code{.bss}.

\textbf{Hints.}
\begin{itemize}
\item One ``flash'' array holds the initial values (LMA); one ``RAM'' array is
  the live \code{.data} (VMA).
\item \code{.bss} is just \code{memset} to 0.
\end{itemize}

\attempt

\textbf{Solution.}
\begin{cgym}
#include <stdio.h>
#include <string.h>
#include <stdint.h>

static uint8_t flash_data_init[4] = { 1, 2, 3, 4 };  /* LMA: init values in flash */
static uint8_t ram_data[4];                           /* VMA: live .data in RAM    */
static uint8_t ram_bss[4]  = { 9, 9, 9, 9 };          /* stand-in for .bss memory  */

static void startup_copy_data(void) {                 /* the .data copy loop */
    memcpy(ram_data, flash_data_init, sizeof ram_data);
}
static void startup_zero_bss(void) {                  /* the .bss zero loop  */
    memset(ram_bss, 0, sizeof ram_bss);
}

int main(void) {
    startup_copy_data();
    startup_zero_bss();
    printf("data=%d%d%d%d bss=%d%d%d%d\n",            /* data=1234 bss=0000 */
           ram_data[0], ram_data[1], ram_data[2], ram_data[3],
           ram_bss[0], ram_bss[1], ram_bss[2], ram_bss[3]);
    return 0;
}
\end{cgym}

On real hardware the bounds of these loops are the linker symbols
(\code{\_sdata}/\code{\_edata}/\code{\_etext} for the copy,
\code{\_sbss}/\code{\_ebss} for the zero), and this runs in
\code{Reset\_Handler} before \code{main}. Comment out
\code{startup\_zero\_bss} and \code{ram\_bss} keeps its old contents
(\code{9999}) --- exactly the ``globals hold garbage from reset'' failure of
Q3.12.

\followups
\begin{enumerate}
\item \emph{``Which linker symbols bound these loops?''} $\rightarrow$
  \code{\_sdata}/\code{\_edata}/\code{\_etext} and \code{\_sbss}/\code{\_ebss}.
\item \emph{``Why copy \code{.data} at all --- why not run from flash?''}
  $\rightarrow$ It's read/write; flash isn't freely writable, so live
  variables must be in RAM.
\item \emph{``What else runs before \code{main} in C++?''} $\rightarrow$
  Static constructors via \code{\_\_libc\_init\_array} / \code{init\_array}.
\end{enumerate}
\end{gymproblem}

%---------------------------------------------------------------
\begin{gymproblem}{O(1) free-list allocator}

\textbf{Problem.} Build a fixed-block allocator whose \code{alloc} and
\code{free} are both O(1) by threading a singly-linked free list through the
blocks themselves.

\textbf{Hints.}
\begin{itemize}
\item Each free block stores a \code{next} pointer (the block memory doubles as
  list storage).
\item \code{alloc} pops the head; \code{free} pushes onto the head.
\end{itemize}

\attempt

\textbf{Solution.}
\begin{cgym}
#include <stdio.h>
#include <stddef.h>

#define N 4
typedef struct block { struct block *next; } block_t;
static block_t  storage[N];
static block_t *free_list;

static void fl_init(void) {
    free_list = NULL;
    for (int i = 0; i < N; i++) {        /* push every block onto the list */
        storage[i].next = free_list;
        free_list = &storage[i];
    }
}
static void *fl_alloc(void) {            /* pop the head -- O(1) */
    if (!free_list) return NULL;
    block_t *b = free_list;
    free_list = b->next;
    return b;
}
static void fl_free(void *p) {           /* push onto the head -- O(1) */
    block_t *b = p;
    b->next = free_list;
    free_list = b;
}

int main(void) {
    fl_init();
    void *a = fl_alloc(), *b = fl_alloc();
    fl_free(a);
    void *c = fl_alloc();
    printf("reused=%d b=%d\n", c == a, b != NULL);   /* reused=1 b=1 */
    return 0;
}
\end{cgym}

The trick is that a \emph{free} block holds no user data, so its bytes can store
the \code{next} link --- the free list costs no extra memory. Both operations
are a couple of pointer moves, fully deterministic. This is how RTOS memory
pools and many lock-free allocators work; the same intrusive-list idea recovers
its parent via \code{container\_of} (Chapter 1).

\followups
\begin{enumerate}
\item \emph{``Make it lock-free for ISR use.''} $\rightarrow$ A compare-and-swap
  on the list head (ABA-aware) or a brief interrupt disable.
\item \emph{``Block size must hold a pointer --- implication?''} $\rightarrow$
  Minimum block size is \code{sizeof(void*)}; tiny objects waste space, so size
  pools to your message types.
\item \emph{``Detect freeing a pointer that isn't ours.''} $\rightarrow$
  Range-check \code{p} against \code{storage} bounds and block alignment.
\end{enumerate}
\end{gymproblem}

%---------------------------------------------------------------
\begin{gymproblem}{Stack high-water mark by painting}

\textbf{Problem.} Implement the RTOS ``stack painting'' technique: fill a
stack region with a sentinel, then later report how many bytes were actually
used.

\textbf{Hints.}
\begin{itemize}
\item Paint the whole region with a constant (e.g.\ \code{0xCD}).
\item Usage = region size minus the count of untouched sentinel bytes from the
  far end.
\end{itemize}

\attempt

\textbf{Solution.}
\begin{cgym}
#include <stdio.h>
#include <stdint.h>
#include <string.h>

#define STK   256
#define PAINT 0xCD
static uint8_t fake_stack[STK];

static void   paint(void) { memset(fake_stack, PAINT, STK); }
static size_t high_water(void) {
    for (size_t i = 0; i < STK; i++)        /* scan from the growth limit up */
        if (fake_stack[i] != PAINT)
            return STK - i;                 /* first touched byte -> usage   */
    return 0;
}

int main(void) {
    paint();
    /* simulate the stack being used from the top (high addresses) down */
    fake_stack[STK - 1] = 1;
    fake_stack[STK - 2] = 2;
    fake_stack[STK - 10] = 3;               /* deepest touch */
    printf("high_water=%zu\n", high_water());   /* high_water=10 */
    return 0;
}
\end{cgym}

Because a downward stack is used from the high end, the deepest byte ever
written is the high-water mark; scanning from the low (limit) end, the first
non-sentinel byte tells you how close the worst case came to the guard.
Subtracting from \code{STK} gives bytes used (10 here) and the untouched
remainder is your safety margin. This is literally how
\code{uxTaskGetStackHighWaterMark}-style APIs work.

\followups
\begin{enumerate}
\item \emph{``Why can this under-report?''} $\rightarrow$ If real data happened
  to equal the sentinel, or the true worst case never ran --- drive the genuine
  worst path.
\item \emph{``What sentinel value is safe?''} $\rightarrow$ One unlikely in
  real data (\code{0xCD}/\code{0xA5}); some RTOSes use \code{0xA5A5A5A5}.
\item \emph{``Combine with a hard guard.''} $\rightarrow$ An MPU no-access
  region or \code{MSPLIM} at the limit catches the overflow the instant it
  happens, not after the fact.
\end{enumerate}
\end{gymproblem}

%---------------------------------------------------------------
\begin{gymproblem}{Overflow-safe bounded string copy (\code{strlcpy}-style)}

\textbf{Problem.} Copy a C string into a fixed destination buffer, never
overflowing it, always NUL-terminating, and reporting whether truncation
happened.

\textbf{Hints.}
\begin{itemize}
\item Copy at most \code{dstsize - 1} bytes, then terminate.
\item Return the source length so the caller can detect truncation
  (\code{return >= dstsize}).
\end{itemize}

\attempt

\textbf{Solution.}
\begin{cgym}
#include <stdio.h>
#include <stddef.h>

static size_t safe_copy(char *dst, const char *src, size_t dstsize) {
    size_t i = 0;
    if (dstsize > 0) {
        for (; i + 1 < dstsize && src[i]; i++)   /* leave room for NUL */
            dst[i] = src[i];
        dst[i] = '\0';                           /* always terminate   */
    }
    size_t srclen = i;
    while (src[srclen]) srclen++;                /* measure full source */
    return srclen;                               /* >= dstsize => truncated */
}

int main(void) {
    char buf[8];
    size_t need = safe_copy(buf, "abcdefghij", sizeof buf);
    printf("[%s] need=%zu truncated=%d\n",
           buf, need, need >= sizeof buf);        /* [abcdefg] need=10 truncated=1 */
    return 0;
}
\end{cgym}

Unlike \code{strcpy} (no bound $\rightarrow$ overflow) and \code{strncpy}
(may leave the result un-terminated and pads with NULs), this always
terminates and tells you if the source didn't fit --- so a too-long packet
field or filename truncates safely instead of smashing the next variable. Buffer
overflow of exactly this kind is the root of a huge fraction of embedded memory
corruption.

\followups
\begin{enumerate}
\item \emph{``Why is \code{strncpy} not a safe drop-in?''} $\rightarrow$ It
  doesn't NUL-terminate when the source fills the buffer, and zero-pads the
  rest --- two surprises.
\item \emph{``What does the caller do on truncation?''} $\rightarrow$ Treat
  \code{>= dstsize} as an error (reject the input / log it), not silently
  continue.
\item \emph{``Make it copy into the middle of a buffer safely.''}
  $\rightarrow$ Pass the remaining space \code{dstsize - offset}; never trust
  the source length.
\end{enumerate}
\end{gymproblem}

%---------------------------------------------------------------
\begin{gymproblem}{Static memory budgeting with \code{sizeof}}

\textbf{Problem.} A driver keeps four port contexts. Compute the static RAM
footprint at compile time and explain the number.

\textbf{Hints.}
\begin{itemize}
\item \code{sizeof} the struct, then the array.
\item Watch for padding from the trailing \code{uint16\_t}.
\end{itemize}

\attempt

\textbf{Solution.}
\begin{cgym}
#include <stdio.h>
#include <stdint.h>

typedef struct {
    uint8_t  rx[64];
    uint8_t  tx[64];
    uint16_t flags;
} port_t;

int main(void) {
    port_t ports[4];
    printf("one=%zu all=%zu\n", sizeof(port_t), sizeof ports);  /* one=130 all=520 */
    return 0;
}
\end{cgym}

\code{one = 130}: 64 + 64 + 2, with no padding because \code{flags} (a
\code{uint16\_t}, alignment 2) already lands at an even offset (128) and 130 is
a multiple of the struct's alignment (2). \code{all = 520} = 4 $\times$ 130 of
fixed, statically-known RAM --- the kind of number you put in a memory budget
and check against the linker \code{.map}, instead of discovering it at run time
via the heap. (Reorder or add a wider member and padding can change this ---
revisit Chapter 1's padding rules.)

\followups
\begin{enumerate}
\item \emph{``Add a \code{uint32\_t} field --- predict the new size.''}
  $\rightarrow$ Alignment to 4 may add tail padding; compute it, then confirm
  with \code{sizeof}.
\item \emph{``Where do these 520 bytes show up in the build?''} $\rightarrow$
  In \code{.bss} (uninitialised) of the \code{size}/\code{.map} output ---
  RAM, not flash.
\item \emph{``Why prefer this over \code{malloc(4 * sizeof(port\_t))}?''}
  $\rightarrow$ Known at link time, no fragmentation, no failure path,
  analysable for certification.
\end{enumerate}
\end{gymproblem}

%---------------------------------------------------------------
\begin{gymproblem}{Region allocator with checkpoint / rewind (free-all)}

\textbf{Problem.} Extend the bump allocator with \code{mark} and
\code{release} so a caller can allocate scratch space and free it all in O(1)
--- the per-frame ``scratch arena'' pattern.

\textbf{Hints.}
\begin{itemize}
\item \code{mark} captures the current offset; \code{release} restores it.
\item All allocations after a mark vanish together.
\end{itemize}

\attempt

\textbf{Solution.}
\begin{cgym}
#include <stdio.h>
#include <stdint.h>
#include <stddef.h>

static uint8_t scratch[128];
static size_t  sp = 0;

static size_t arena_mark(void)            { return sp; }     /* checkpoint */
static void   arena_release(size_t mark)  { sp = mark; }     /* free-all to mark */
static void  *scratch_alloc(size_t n) {
    if (sp + n > sizeof scratch) return NULL;
    void *p = &scratch[sp];
    sp += n;
    return p;
}

int main(void) {
    size_t m = arena_mark();
    scratch_alloc(40);
    scratch_alloc(20);
    arena_release(m);                     /* release both allocations at once */
    printf("after_release=%zu\n", arena_mark());   /* after_release=0 */
    return 0;
}
\end{cgym}

This is region-based memory management: you allocate freely within a phase
(parsing a packet, building one control-loop frame), then drop the whole region
in a single assignment. No per-object free, no fragmentation, no leaks possible
within the region --- ideal for transient scratch in a real-time loop where
\code{malloc}/\code{free} are banned. The cost is discipline: nothing allocated
after a \code{mark} may outlive the matching \code{release}.

\followups
\begin{enumerate}
\item \emph{``What breaks if a pointer escapes past \code{release}?''}
  $\rightarrow$ It dangles --- the bytes get reused by the next allocation;
  region allocators trade safety for speed.
\item \emph{``Nest marks?''} $\rightarrow$ Yes --- marks form a stack
  (LIFO), so inner regions release before outer.
\item \emph{``Compare to the pool allocator.''} $\rightarrow$ Pool = many
  same-size objects with individual free; arena = mixed sizes, bulk free.
\end{enumerate}
\end{gymproblem}

%=====================================================================
\section{Link to Her Resume}
%=====================================================================

\begin{resumelink}
\textbf{Startup / board bring-up:} ``Board bring-up is exactly where the
memory map stops being theory --- the difference between code that runs under
the debugger and code that runs from reset is whether the startup code copied
\code{.data} and zeroed \code{.bss} correctly. That's the first thing I'd check
on a board that boots in the lab rig but not standalone.''

\medskip
\textbf{Static memory / no-\code{malloc} discipline:} ``Flight-grade firmware
avoids the heap, so buffers come from static pools and arenas sized at link
time. I can read a \code{.map} file and a \code{size} output to confirm flash
and RAM budgets rather than discovering limits at run time.''

\medskip
\textbf{Corruption diagnosis:} ``The hard timing and integration defects I
worked on at BEL are the same family as stack-overflow and buffer-overrun bugs
--- the corruption shows up far from its cause. I reach for guard words, stack
painting, and the \code{.map} file to localise them.''

\medskip
\small Maps to concrete resume lines: TASL board bring-up across autopilot
generations, BEL HSI timing-defect resolution, and the deterministic,
analysable footprint that safety-critical work demands. Keep claims to what the
resume states; the deep project scripting is Chapter 15, from
\code{inputs/INTAKE.md} only.
\end{resumelink}

%=====================================================================
\section{Self-Test Scorecard}
%=====================================================================

\begin{scorecard}
\begin{enumerate}
\item Name the six places data lives in a running C program and one example of
  each.
\item Which segment costs flash but not RAM, and which costs RAM but not flash?
\item List the four build stages and the tool for each.
\item What three things does the startup code do before \code{main}?
\item Define VMA vs LMA using \code{.data}.
\item What does a linker script define, and name two boundary symbols it
  exports.
\item ``Undefined reference'' vs ``multiple definition'' --- which stage, and a
  cause of each.
\item A global is 0 under the debugger but garbage from reset --- what's wrong?
\item Name two ways to detect a stack overflow on an MMU-less target.
\item Why is a fixed-block pool or arena preferred over \code{malloc} in flight
  software?
\end{enumerate}
\end{scorecard}

\begin{answerkey}
\begin{enumerate}
\item \code{.text} (code), \code{.rodata} (\code{const}), \code{.data}
  (\code{int g=5}), \code{.bss} (\code{int g;}), stack (locals), heap
  (\code{malloc}).
\item \code{.bss} costs RAM but not flash (it's just zeroed); the \code{.data}
  \emph{initialisers} cost flash, and \code{.data} costs RAM too. (Pure code
  costs flash but not RAM.)
\item Preprocess (\code{cpp}), compile (\code{cc1}), assemble (\code{as}),
  link (\code{ld}).
\item Set the stack pointer, copy \code{.data} flash$\rightarrow$RAM, zero
  \code{.bss} (then run constructors, call \code{main}).
\item VMA = run address (where \code{.data} lives in RAM); LMA = load address
  (where its init values sit in flash); startup copies LMA$\rightarrow$VMA.
\item The memory regions and section placement (addresses/sizes); exports e.g.\
  \code{\_sdata}, \code{\_edata}, \code{\_sbss}, \code{\_ebss}, \code{\_etext}.
\item Both are \emph{link}-stage. Undefined reference: symbol used but never
  defined/linked. Multiple definition: same external symbol defined in two
  translation units.
\item The startup \code{.bss} zeroing is missing/misconfigured (or its linker
  symbols are wrong); the debugger cleared RAM, reset did not.
\item Any two of: stack canary/guard word, stack painting + high-water scan,
  MPU guard region, \code{MSPLIM}/\code{PSPLIM} stack-limit registers.
\item Deterministic timing, no fragmentation, statically-known and
  link-time-analysable footprint, and no in-flight allocation-failure path ---
  required for certification.
\end{enumerate}
\end{answerkey}

\begin{partnote}
\emph{End of Chapter 2.} Next: \textbf{ARM \& MCU Architecture} --- the
Cortex-M register set, exception model and NVIC, clock tree, GPIO internals,
and the LPC3250/ARM9 specifics behind the BEL platform, where this memory map
is wired to real silicon.
\end{partnote}
